\documentclass[11pt,a4paper]{article}

\usepackage[T1]{fontenc}
\usepackage{lmodern}
\usepackage[a4paper,margin=25mm,headheight=15pt]{geometry}
\usepackage{amsmath,amssymb,amsthm,mathtools,bm}
\usepackage{booktabs,array,longtable}
\usepackage{enumitem}
\usepackage{microtype}
\usepackage{xcolor}
\usepackage{hyperref}
\usepackage{fancyhdr}

\hypersetup{
  colorlinks=true,
  linkcolor=blue!55!black,
  citecolor=blue!55!black,
  urlcolor=blue!55!black,
  pdftitle={An Unconditional Proof of the Logarithmic Morris Constant-Term Identity},
  pdfauthor={Yongcheng Hu},
  pdfsubject={An unconditional proof with a semantic audit of the Lean 4.28 formalization},
  pdfkeywords={constant term, Morris identity, Dyson identity, Pfaffian, Lean}
}

\pagestyle{fancy}
\fancyhf{}
\fancyhead[L]{The Logarithmic Morris Constant-Term Identity}
\fancyhead[R]{An Unconditional Proof}
\fancyfoot[C]{\thepage}

\setlength{\parindent}{1.5em}
\setlength{\parskip}{0.22em}
\setlist{nosep,leftmargin=2.2em}
\allowdisplaybreaks[3]

\newtheorem{theorem}{Theorem}[section]
\newtheorem{lemma}[theorem]{Lemma}
\newtheorem{proposition}[theorem]{Proposition}
\newtheorem{corollary}[theorem]{Corollary}
\theoremstyle{definition}
\newtheorem{definition}[theorem]{Definition}
\newtheorem{remark}[theorem]{Remark}

\newcommand{\CT}{\operatorname{CT}}
\newcommand{\Alt}{\operatorname{Alt}}
\newcommand{\Pf}{\operatorname{Pf}}
\newcommand{\supp}{\operatorname{supp}}
\newcommand{\sgn}{\operatorname{sgn}}
\newcommand{\e}{\mathrm e}
\newcommand{\T}{\mathbb T}
\newcommand{\Z}{\mathbb Z}
\newcommand{\Q}{\mathbb Q}
\newcommand{\C}{\mathbb C}
\newcommand{\E}{\mathbb E}
\newcommand{\dd}{\,\mathrm d}
\newcommand{\Logc}{\operatorname{Log}_{\circ}}

\title{\bfseries An Unconditional Proof of the Logarithmic Morris
Constant-Term Identity\\[5pt]
\large Finite Laurent Coefficients, Augmented Pfaffians,
and Compactly Supported Moments}
\author{Yongcheng Hu\\
Central South University\\
\href{mailto:252111040@csu.edu.cn}{\texttt{252111040@csu.edu.cn}}}
\date{Version of 22 July 2026}

\begin{document}
\maketitle

\begin{abstract}
Let $n=2m+1$ and $K=2k+1$, where $a,b,k,m$ are nonnegative
integers. We give an unconditional proof of the logarithmic Morris
constant-term identity. The proof uses neither the complex Morris
constant-term conjecture nor the evaluation of the Dyson integral for
arbitrary complex parameters.

We first define the constant term involving the formal logarithmic product
as a linear functional on the finite support of a Laurent polynomial. We
then compute the Fourier coefficients of the one-variable principal
logarithm and justify the boundary limit by an explicit uniform
$L^1$ majorant. The central evaluation of the alternating logarithmic sum
uses the Pfaffian of an augmented skew-symmetric matrix: a congruence
transformation of determinant $1$ eliminates all angular variables and
avoids an informal combinatorial involution. Circularization is performed
pointwise only on the full-measure open set of pairwise distinct
coordinates, with integrability proved separately.

In the Dyson specialization $a=b=0$, Good's finite induction gives all
integer moments of the circular discriminant. We construct a product of
mutually independent Beta random variables with the same moments. After
pushing both distributions forward to a common compact interval, the Weierstrass
approximation theorem yields moment determinacy, allowing the integer
moment formula to be used at the required half-integer exponent. Finally,
Laurent Euler derivatives, a finite weight-shift identity for the logarithmic
weight, and a three-variable permutation invariance after omitting one
logarithm produce a recurrence in the endpoint parameters. A finite
Chu--Vandermonde identity and symmetry in $a,b$ complete the proof.

The final section records the interface-level correspondence between the
mathematical objects used here and their Lean definitions. The project was
built under \path{leanprover/lean4:v4.28.0}; the final theorem contains no
\texttt{sorry}, and its axiom report contains only three standard logical
dependencies of Lean/Mathlib.
\end{abstract}

\tableofcontents

\section{Main Theorem, Conventions, and the Rank-One Case}

Let $X=(x_1,\ldots,x_n)$ and set
\[
  \Delta(X)=\prod_{1\le i<j\le n}(x_i-x_j).
\]
Throughout the paper,
\[
  n=2m+1,\qquad K=2k+1,\qquad
  a,b,k,m\in\mathbb Z_{\ge 0}.
\]
We use the double-factorial conventions
\[
  0!!=1!!=1,\qquad (N+2)!!=(N+2)N!!.
\]
Define the finite Laurent polynomial
\begin{equation}\label{eq:kernel}
\Phi_{a,b,k}(X)
=\Delta(X)\prod_{i=1}^{n}x_i^{-m}(1-x_i)^a(1-x_i^{-1})^b
 \prod_{\substack{1\le i,j\le n\\i\ne j}}
 \left(1-\frac{x_i}{x_j}\right)^k .
\end{equation}

\begin{theorem}[Logarithmic Morris constant-term identity]\label{thm:main}
For the parameters above,
\begin{align}
&\CT\!\left[
 \Phi_{a,b,k}(X)
 \prod_{q=1}^{m}\log\!\left(1-\frac{x_{2q}}{x_{2q-1}}\right)
\right]\notag\\
&\qquad=
\frac1{n!!}\prod_{i=0}^{n-1}
\frac{(2a+2b+iK)!!\,((i+1)K)!!}
{(2a+iK)!!\,(2b+iK)!!\,K!!}.
\label{eq:main}
\end{align}
\end{theorem}

\noindent\textbf{Relation to the original theorem.}
Theorem~\ref{thm:main} is precisely the assertion of
\cite[Theorem~1.3]{LogMorris} after writing $n=2m+1$ and $K=2k+1$.
In contrast to that theorem's original derivation, the present proof does
not assume the complex Morris constant-term identity
\cite[(4.5)]{LogMorris}.

\begin{remark}[Meaning of the formal logarithm]\label{rem:formal-log}
On the left-hand side of \eqref{eq:main}, the logarithm is understood as
\[
  \log(1-z)=-\sum_{r\ge1}\frac{z^r}{r}.
\]
This does not mean that one first chooses a multivariate formal power
series ring and then performs an infinite summation. Since
$\Phi_{a,b,k}$ has finite Laurent support, only finitely many exponents can
contribute to the constant term. Section~\ref{sec:finite-functional}
defines it directly as a finite sum.
\end{remark}

\begin{proof}[The rank-one case $n=1$]
Here $m=0$, so there are no logarithmic, Vandermonde, or interaction
factors. Hence the left-hand side is
\[
 \CT\bigl[(1-x)^a(1-x^{-1})^b\bigr]
 =\sum_{r\ge0}\binom ar\binom br
 =\binom{a+b}{a},
\]
where the last equality is the finite Vandermonde convolution. The
right-hand side has only the term $i=0$ and equals
\[
 \frac{(2a+2b)!!}{(2a)!!(2b)!!}
 =\frac{2^{a+b}(a+b)!}{2^aa!\,2^bb!}
 =\binom{a+b}{a}.
\]
Thus the theorem holds for $n=1$. Unless explicitly stated otherwise, we
assume henceforth that $m\ge1$, and therefore $n\ge3$.
\end{proof}

\section{The Finite Logarithmic Constant-Term Functional}
\label{sec:finite-functional}

Let
\[
 F(X)=\sum_{\alpha\in\Z^n}c_\alpha X^\alpha,\qquad
 X^\alpha=\prod_{i=1}^n x_i^{\alpha_i},
\]
where only finitely many $c_\alpha$ are nonzero. Define
\begin{equation}\label{eq:weight}
w(\alpha)=
\begin{cases}
\displaystyle\prod_{q=1}^{m}\frac{-1}{\alpha_{2q-1}},
 &\begin{array}{l}
 \alpha_{2q-1}>0,\quad
 \alpha_{2q}=-\alpha_{2q-1}\quad(1\le q\le m),\\[-2pt]
 \alpha_n=0,
 \end{array}\\[9pt]
0,&\text{otherwise},
\end{cases}
\end{equation}
and set
\begin{equation}\label{eq:Lambda}
 \Lambda(F)=\sum_{\alpha\in\supp(F)}c_\alpha w(\alpha).
\end{equation}

\begin{lemma}\label{lem:formal-ct}
For every finite Laurent polynomial $F$,
\[
\Lambda(F)=
\CT\!\left[
F(X)\prod_{q=1}^{m}
\log\!\left(1-\frac{x_{2q}}{x_{2q-1}}\right)
\right].
\]
\end{lemma}

\begin{proof}
The $q$th logarithm contributes precisely the monomials
\[
 -\frac{x_{2q}^{r_q}x_{2q-1}^{-r_q}}{r_q},
 \qquad r_q>0.
\]
Its product with $X^\alpha$ has exponent zero in coordinates
$2q-1,2q$ if and only if
\[
 \alpha_{2q-1}=r_q>0,\qquad
 \alpha_{2q}=-r_q.
\]
The unpaired coordinate requires $\alpha_n=0$. Under these conditions
the coefficient is exactly
$\prod_q(-1/\alpha_{2q-1})$; otherwise the contribution is zero.
Because $\supp(F)$ is finite, each $\alpha$ corresponds to at most one
tuple $(r_1,\ldots,r_m)$. Thus the formal constant term is precisely the
finite sum \eqref{eq:Lambda}.
\end{proof}

We henceforth write
\[
  L(a,b):=\Lambda(\Phi_{a,b,k}).
\]

\section{The Fourier Bridge and Finite Coefficient Extraction}
\label{sec:fourier}

Let $\T=\mathbb R/\mathbb Z$, equipped with normalized Haar measure
$\dd t$, and write
\[
  z(t)=\e^{2\pi i t}.
\]
Let $\operatorname{Log}$ denote the principal complex logarithm. To obtain
an everywhere-defined measurable function, put
\[
\Logc(\zeta)=
\begin{cases}
\operatorname{Log}(\zeta),&\zeta\ne0,\\
0,&\zeta=0.
\end{cases}
\]
The value at zero does not change any integral; it also agrees with the
total-function convention used by Lean's \texttt{Complex.log}. Set
\[
  \ell(t)=\Logc(1-z(t)).
\]

\begin{lemma}[$L^1$ integrability and Fourier coefficients]
\label{lem:fourier}
The function $\ell$ belongs to $L^1(\T)$ and, for every $q\in\Z$,
\begin{equation}\label{eq:fourier}
\int_\T z(t)^q\ell(t)\dd t=
\begin{cases}
1/q,&q<0,\\
0,&q\ge0.
\end{cases}
\end{equation}
\end{lemma}

\begin{proof}
We first prove integrability. For $t\ne0$,
\[
 |1-\e^{2\pi it}|=2|\sin(\pi t)|.
\]
In a representative neighborhood of $t=0$, the quantities
$|\sin(\pi t)|$ and $|t|$ are comparable. Consequently,
\[
 |\ell(t)|\le C+|\log|t||.
\]
Since $\int_0^\varepsilon|\log t|\dd t<\infty$, while $\ell$ is
continuous and bounded away from $0$, this proves $\ell\in L^1(\T)$.

For $0\le r<1$, let
\[
  \ell_r(t)=\operatorname{Log}(1-rz(t)).
\]
The absolutely convergent expansion inside the unit disk gives
\[
  \ell_r(t)=-\sum_{j\ge1}\frac{r^jz(t)^j}{j}.
\]
Termwise integration and character orthogonality yield
\begin{equation}\label{eq:radial-fourier}
\int_\T z(t)^q\ell_r(t)\dd t=
\begin{cases}
-r^{-q}/(-q),&q<0,\\
0,&q\ge0.
\end{cases}
\end{equation}

It remains to justify the limit $r\uparrow1$. Fix $t\ne0$ and set
\[
 d=|1-z(t)|,\qquad u=|1-rz(t)|.
\]
The reverse triangle inequality gives $1-r\le u$, and
\[
 d\le |1-rz(t)|+|(r-1)z(t)|\le2u,
\qquad u\le1+r\le2.
\]
Hence
\[
 |\log u|\le\log2+|\log d|.
\]
The absolute value of the principal argument is at most $\pi$, so
\begin{equation}\label{eq:domination}
|\ell_r(t)|
\le |\log u|+\pi
\le \log2+|\ell(t)|+\pi.
\end{equation}
The right-hand side belongs to $L^1(\T)$ and is independent of $r$.
Moreover, $\ell_r(t)\to\ell(t)$ away from zero. Dominated convergence
together with \eqref{eq:radial-fourier} now gives \eqref{eq:fourier}.
\end{proof}

\begin{lemma}[Two-variable Fourier coefficient]\label{lem:pair-fourier}
For $p,q\in\Z$,
\begin{equation}\label{eq:pair-fourier}
\int_{\T^2}z(u)^pz(v)^q
\Logc\!\left(1-\frac{z(v)}{z(u)}\right)\dd u\dd v
=
\begin{cases}
-1/p,&p>0,\ q=-p,\\
0,&\text{otherwise}.
\end{cases}
\end{equation}
\end{lemma}

\begin{proof}
Apply the Haar-measure-preserving shear
$(u,s)\mapsto(u,u+s)$. The integrand becomes
\[
 z(u)^{p+q}\,z(s)^q\ell(s).
\]
Fubini's theorem separates the two variables. The first integral forces
$p+q=0$, while the second is evaluated by Lemma~\ref{lem:fourier}.
When $q=-p$, it is nonzero precisely for $p>0$, and its value is
$1/q=-1/p$.
\end{proof}

\begin{proposition}[Finite-coefficient circular-integral bridge]
\label{prop:bridge}
Define
\[
 \mathcal L(t)=\prod_{q=1}^{m}
 \Logc\!\left(1-\frac{z(t_{2q})}{z(t_{2q-1})}\right).
\]
Then, for every finite Laurent polynomial $F$,
\begin{equation}\label{eq:bridge}
 \Lambda(F)=\int_{\T^n}F(z(t_1),\ldots,z(t_n))
 \mathcal L(t)\dd t.
\end{equation}
\end{proposition}

\begin{proof}
The pairs of variables are disjoint. Put
\[
 f_q(t_{2q-1},t_{2q})
 =\Logc\!\left(1-\frac{z(t_{2q})}{z(t_{2q-1})}\right).
\]
Lemma~\ref{lem:fourier} and the measure-preserving shear show that each
$f_q\in L^1(\T^2)$. By Tonelli's theorem and the disjointness of the
variable blocks,
\[
 \int_{\T^{2m}}\prod_{q=1}^m|f_q|
 =\prod_{q=1}^m\int_{\T^2}|f_q|<\infty.
\]
The unpaired coordinate contributes the factor $1$, so
$\mathcal L\in L^1(\T^n)$; Fubini's theorem may therefore be used to
separate the variable blocks. Since a Laurent polynomial is bounded on
the compact torus, the integrand in \eqref{eq:bridge} is integrable.

First take $F=X^\alpha$. Apply \eqref{eq:pair-fourier} to each pair and
character orthogonality to the unpaired variable. The integral is nonzero
exactly under the admissibility conditions in \eqref{eq:weight}, and its
value is $w(\alpha)$. The general case follows by finite linearity.
\end{proof}

\section{Augmented Pfaffians and Circularization}
\label{sec:circularization}

For an arbitrary $n\times n$ matrix $A=(A_{ij})$, define
\begin{equation}\label{eq:alt}
\Alt_m(A)=\sum_{\sigma\in S_n}\sgn(\sigma)
\prod_{q=1}^{m}A_{\sigma(2q-1),\sigma(2q)}.
\end{equation}
For $c,y_1,\ldots,y_n\in\C$, define the skew-symmetric matrix
\begin{equation}\label{eq:saw}
S_{ij}(c,y)=
\begin{cases}
y_j-y_i+c,&i<j,\\
-(y_i-y_j+c),&j<i,\\
0,&i=j.
\end{cases}
\end{equation}

\begin{lemma}[Augmented sawtooth Pfaffian]\label{lem:saw}
If $n=2m+1$, then
\begin{equation}\label{eq:saw-eval}
 \Alt_m(S(c,y))=2^m m!c^m.
\end{equation}
\end{lemma}

\begin{proof}
Introduce the index $N=n+1=2m+2$ and define the $N\times N$
skew-symmetric matrix $B$ by
\[
 B_{ij}=S_{ij}(c,y)\quad(1\le i,j\le n),\qquad
 B_{iN}=1,\quad B_{Ni}=-1.
\]
Every Pfaffian principal submatrix below is ordered by the inherited
natural order of its indices, and we set $\Pf(\varnothing)=1$. We use the
standard convention
\[
\Pf(B)=\frac1{2^{m+1}(m+1)!}
\sum_{\tau\in S_N}\sgn(\tau)
\prod_{r=1}^{m+1}B_{\tau(2r-1),\tau(2r)}.
\]
Group \eqref{eq:alt} according to the unpaired vertex
$u=\sigma(n)$. Moving $u$ from its $u$th position in the natural order
to the last position requires $n-u$ adjacent transpositions. Since $n$
is odd,
\[
(-1)^{n-u}=(-1)^{u+1}.
\]
For fixed $u$, the sum over oriented, ordered pairings of the remaining
$2m$ vertices is
$2^m m!\Pf(B_{\widehat u,\widehat N})$, where
$B_{\widehat u,\widehat N}$ denotes the principal submatrix obtained by
deleting rows and columns $u,N$. Therefore
\[
\Alt_m(S)=2^m m!\sum_{u=1}^{n}
(-1)^{u+1}\Pf(B_{\widehat u,\widehat N}).
\]
On the other hand, expansion of the augmented matrix $B$ along its last
row and column, using $B_{uN}=1$, gives
\[
\Pf(B)=\sum_{u=1}^{n}(-1)^{u+1}
\Pf(B_{\widehat u,\widehat N}).
\]
The signs agree term by term, so
\begin{equation}\label{eq:alt-pf}
 \Alt_m(S)=2^m m!\Pf(B).
\end{equation}

Regard $B$ as its associated alternating bilinear form. Let
$e_1,\ldots,e_N$ be the standard basis and make the change of basis
\[
 f_i=e_i-y_i e_N\quad(1\le i\le n),\qquad f_N=e_N.
\]
Its transition matrix is unitriangular and has determinant $1$. The
Pfaffian congruence formula
\[
 \Pf(P^{\mathsf T}BP)=\det(P)\Pf(B)
\]
therefore leaves the Pfaffian unchanged. Standard conventions for
Pfaffian expansion and congruence may be found in \cite{Knuth1996}.
For $i<j\le n$, a direct calculation gives
\[
\begin{aligned}
B(f_i,f_j)
&=B(e_i,e_j)-y_jB(e_i,e_N)-y_iB(e_N,e_j)\\
&=(y_j-y_i+c)-y_j+y_i=c,
\end{aligned}
\]
and $B(f_i,f_N)=1$. Thus the strictly upper-triangular part of the new
matrix $C$ satisfies
\[
 C_{ij}=c\quad(1\le i<j\le n),\qquad C_{iN}=1.
\]
Every perfect matching contains exactly one edge incident to $N$; the
remaining $m$ edges lie among the first $n$ vertices. Consequently every
Pfaffian term contains the factor $c^m$, and
\[
 \Pf(C)=c^m\Pf(J_N),
\]
where $J_N$ is the skew-symmetric matrix whose strictly upper-triangular
entries are all $1$.

Finally, we show that $\Pf(J_{2r})=1$. For $r=0$, this is precisely the
empty-Pfaffian convention. Expanding along the first row and using the
induction hypothesis,
\[
 \Pf(J_{2r})
 =\sum_{j=2}^{2r}(-1)^j\Pf(J_{2r-2})
 =\sum_{j=2}^{2r}(-1)^j=1.
\]
Hence $\Pf(B)=\Pf(C)=c^m$. Substitution into \eqref{eq:alt-pf} proves
\eqref{eq:saw-eval}.
\end{proof}

Define the circular kernel
\begin{equation}\label{eq:circular-kernel}
\mathcal C_{a,b,K}(t)=
\prod_{i=1}^{n}(1-z(t_i))^a(1-z(t_i)^{-1})^b
\prod_{1\le i<j\le n}|z(t_i)-z(t_j)|^K.
\end{equation}
When $a\ne b$, it is generally complex-valued; for this reason we call it
a circular kernel rather than a density.

\begin{lemma}[The alternating logarithmic sum in an ordered chamber]
\label{lem:ordered-alt}
Suppose
\[
 0\le\theta_1<\cdots<\theta_n<2\pi,\qquad z_j=\e^{i\theta_j},
\]
and put
\[
 A_{ij}=\Logc(1-z_j/z_i)\qquad(1\le i,j\le n).
\]
In particular, the total-function convention of
Section~\ref{sec:fourier} gives $A_{ii}=\Logc(0)=0$. Then
\begin{equation}\label{eq:log-alt}
 \Alt_m(A)=m!(-\pi i)^m.
\end{equation}
\end{lemma}

\begin{proof}
For $i<j$, let
\[
 \delta=\theta_j-\theta_i\in(0,2\pi).
\]
We have
\[
\begin{aligned}
1-\e^{i\delta}
 &=2\sin\frac{\delta}{2}\,
   \e^{\,i(\delta/2-\pi/2)},\\
1-\e^{-i\delta}
 &=2\sin\frac{\delta}{2}\,
   \e^{\,i(\pi/2-\delta/2)}.
\end{aligned}
\]
Both displayed arguments lie in $(-\pi,\pi)$, so the principal branch
gives
\[
 \operatorname{Log}(1-\e^{i\delta})
 -\operatorname{Log}(1-\e^{-i\delta})
 =i(\delta-\pi).
\]
Both arguments of the logarithms are nonzero, so $\Logc$ agrees here
with the principal logarithm. Therefore
\begin{equation}\label{eq:log-skew}
 A_{ij}-A_{ji}=i(\theta_j-\theta_i-\pi).
\end{equation}
Independently swapping the two positions in every input pair of
\eqref{eq:alt} gives
\[
 \Alt_m(A-A^{\mathsf T})=2^m\Alt_m(A).
\]
By \eqref{eq:log-skew},
\[
A-A^{\mathsf T}=S(-\pi i,(i\theta_j)_j).
\]
Lemma~\ref{lem:saw}, followed by cancellation of $2^m$, proves
\eqref{eq:log-alt}.
\end{proof}

\begin{lemma}[The phase in an ordered chamber]\label{lem:ordered-phase}
In the ordered chamber of Lemma~\ref{lem:ordered-alt},
\begin{equation}\label{eq:kernel-phase}
 \Phi_{a,b,k}(z_1,\ldots,z_n)
 =(-i)^{mn}\mathcal C_{a,b,K}(t).
\end{equation}
\end{lemma}

\begin{proof}
For each $i<j$,
\[
\left(1-\frac{x_i}{x_j}\right)^k
\left(1-\frac{x_j}{x_i}\right)^k
=(-1)^k\frac{(x_i-x_j)^{2k}}{(x_ix_j)^k}.
\]
Hence
\begin{equation}\label{eq:normalized-kernel}
\Phi_{a,b,k}
=\varepsilon\,\Delta^K
\prod_{i=1}^{n}x_i^{-(b+mK)}(1-x_i)^{a+b},
\qquad
\varepsilon=(-1)^{k\binom n2+bn}.
\end{equation}
In the ordered chamber, for every pair we use
\[
 z_i-z_j=(-i)\e^{i(\theta_i+\theta_j)/2}|z_i-z_j|.
\]
In the midpoint phase of $\Delta^K$, each $z_i$ occurs
$n-1=2m$ times, exactly cancelling $z_i^{-mK}$. Rewriting
$(1-z_i)^{a+b}z_i^{-b}$ as
$(-1)^b(1-z_i)^a(1-z_i^{-1})^b$, and setting
$C=\binom n2$, we find that the endpoint factors contribute
$(-1)^{bn}$ and the pairwise decomposition of $\Delta^K$ contributes
$(-i)^{KC}$. Together with $\varepsilon$ from
\eqref{eq:normalized-kernel}, the total phase is
\[
\begin{aligned}
\varepsilon(-1)^{bn}(-i)^{KC}
 &=(-1)^{kC}(-i)^{KC}\\
 &=(-i)^{2kC}(-i)^{(2k+1)C}\\
 &=(-i)^{(4k+1)C}\\
 &=(-i)^C.
\end{aligned}
\]
Since $C=\binom n2=mn$, this is $(-i)^{mn}$, proving
\eqref{eq:kernel-phase}.
\end{proof}

\begin{proposition}[Circularization formula]\label{prop:circularization}
For all $a,b,k,m\in\mathbb Z_{\ge0}$,
\begin{equation}\label{eq:circularization}
 L(a,b)=\frac{m!\pi^m}{n!}
 \int_{\T^n}\mathcal C_{a,b,K}(t)\dd t.
\end{equation}
\end{proposition}

\begin{proof}
For $\sigma\in S_n$ and $t\in\T^n$, fix the convention
\[
 (\sigma t)_j:=t_{\sigma(j)}\qquad(1\le j\le n).
\]
By Proposition~\ref{prop:bridge}, $L(a,b)$ is the integral of
$\Phi(z(t))\mathcal L(t)$. The Laurent kernel is bounded on the compact
torus, and $\mathcal L$ is integrable; hence their product is integrable.
Define
\[
\mathcal A(t)=\sum_{\sigma\in S_n}\sgn(\sigma)\,
\mathcal L(\sigma t).
\]
Changes of variables by permutations preserve Haar measure. Since
$\Phi$ is alternating, every
$\Phi(t)\mathcal L(\sigma t)$ is integrable, and the finite alternating
sum may be integrated term by term:
\begin{equation}\label{eq:average}
n!\int_{\T^n}\Phi(z(t))\mathcal L(t)\dd t
=\int_{\T^n}\Phi(z(t))\mathcal A(t)\dd t.
\end{equation}

In any ordered chamber, let
\[
 A_{ij}=\Logc\!\left(1-\frac{z(t_j)}{z(t_i)}\right).
\]
Termwise comparison of the definitions of $\mathcal L$, $\mathcal A$,
and \eqref{eq:alt} gives
\begin{equation}\label{eq:A-is-alt}
 \mathcal A(t)=\Alt_m(A).
\end{equation}
This is the precise interface between alternation under the integral and
Lemma~\ref{lem:ordered-alt}.

Let $U\subset\T^n$ be the set of tuples with pairwise distinct
coordinates. Its complement is a finite union of diagonal sets
$\{t_i=t_j\}$, each of Haar measure zero, so $U$ has full measure. For
$t\in U$, choose the unique representative
$\theta_i\in[0,2\pi)$ of each $t_i\in\T=\mathbb R/\mathbb Z$, with
$t_i=\theta_i/(2\pi)$, and then the unique sorting permutation $\rho$
for which the representatives of
$(\rho t)_1,\ldots,(\rho t)_n$ are strictly increasing. By the convention
above and by alternation,
\[
\Phi(\rho t)=\sgn(\rho)\Phi(t),\qquad
\mathcal A(\rho t)=\sgn(\rho)\mathcal A(t),
\]
whereas
$\mathcal C_{a,b,K}(\rho t)=\mathcal C_{a,b,K}(t)$. Thus
$\Phi\mathcal A$ is permutation-invariant, and
Lemmas~\ref{lem:ordered-alt} and \ref{lem:ordered-phase} apply in the
sorted chamber. Since
\[
(-i)^{mn}(-i)^m=(-i)^{m(n+1)}=1
\]
and $m(n+1)=2m(m+1)$ is divisible by $4$, we obtain on $U$
\[
 \Phi(z(t))\mathcal A(t)=m!\pi^m\mathcal C_{a,b,K}(t).
\]
The function $\mathcal C_{a,b,K}$ is Borel measurable, since it is a
finite product of continuous functions on $\T^n$.
This is an almost-everywhere identity; no meaningless product
$0\cdot\infty$ is asserted at collision points. The left-hand side is
integrable, so the right-hand side is integrable as well. Substitution
into \eqref{eq:average}, followed by Proposition~\ref{prop:bridge},
gives \eqref{eq:circularization}.
\end{proof}

\begin{corollary}[Exchange of endpoint parameters]\label{cor:swap}
\[
 L(a,b)=L(b,a).
\]
\end{corollary}

\begin{proof}
In \eqref{eq:circularization}, make the Haar-measure-preserving change of
variables $t\mapsto-t$. Then $z_i\mapsto z_i^{-1}$, the absolute
Vandermonde factor is unchanged, and the two endpoint factors are
interchanged.
\end{proof}

\section{An Unconditional Evaluation in the Dyson Specialization}
\label{sec:dyson}

In this section set $a=b=0$ and define
\begin{equation}\label{eq:Y}
Y(t)=\prod_{1\le i<j\le n}|z(t_i)-z(t_j)|^2.
\end{equation}
Then $\mathcal C_{0,0,K}(t)=Y(t)^{K/2}$. We begin only with the finite
Dyson identity at nonnegative integer exponents; no evaluation at general
complex parameters is assumed.

\subsection{Good's Finite Dyson Identity}

For a finite index set $I$ and
$\bm a=(a_i)_{i\in I}\in\mathbb Z_{\ge0}^{\,I}$, put
\[
 R_i(X)=\prod_{\substack{j\in I\\j\ne i}}
 \left(1-\frac{x_i}{x_j}\right),\qquad
 D_{I,\bm a}(X)=\prod_{i\in I}R_i(X)^{a_i}.
\]

\begin{proposition}[Good--Dyson]\label{prop:good}
\begin{equation}\label{eq:good}
 \CT[D_{I,\bm a}]
 =\frac{\left(\sum_{i\in I}a_i\right)!}{\prod_{i\in I}a_i!}.
\end{equation}
\end{proposition}

\begin{proof}
If $I=\varnothing$, both sides are $1$. Assume henceforth that
$I\ne\varnothing$ and regard the $x_i$ as algebraically independent.
The sum of the Lagrange basis polynomials evaluated at $0$ gives, first
in the rational function field $\Q(x_i\mid i\in I)$,
\[
 1=\sum_{i\in I}\frac1{R_i}.
\]
Clearing the denominator $\prod_{i\in I}R_i$ turns this into the Laurent
polynomial identity
\[
 \prod_{i\in I}R_i
 =\sum_{i\in I}\prod_{\substack{j\in I\\j\ne i}}R_j.
\]
If every $a_i>0$, multiplication by $\prod_iR_i^{a_i-1}$ yields
\begin{equation}\label{eq:good-rec}
 D_{I,\bm a}=\sum_{i\in I}D_{I,\bm a-\bm e_i}.
\end{equation}
If $a_h=0$, then
\[
D_{I,\bm a}
=D_{I\setminus\{h\},\,\bm a|_{I\setminus\{h\}}}
\prod_{i\in I\setminus\{h\}}
\left(1-\frac{x_i}{x_h}\right)^{a_i}.
\]
The first factor is independent of $x_h$, while the second contains only
nonpositive powers of $x_h$ and has constant term $1$ in $x_h$. Hence
\begin{equation}\label{eq:good-boundary}
\CT[D_{I,\bm a}]
=\CT[D_{I\setminus\{h\},\,\bm a|_{I\setminus\{h\}}}].
\end{equation}

The candidate
\[
G(I,\bm a)=\frac{(\sum_i a_i)!}{\prod_i a_i!}
\]
satisfies the same recurrence and boundary relation. Indeed, if every
$a_i>0$ and $A=\sum_i a_i$, then
\[
\sum_iG(I,\bm a-\bm e_i)
=G(I,\bm a)\sum_i\frac{a_i}{A}=G(I,\bm a).
\]
If $a_h=0$, then $0!=1$ gives directly
\[
 G(I,\bm a)
 =G(I\setminus\{h\},\,\bm a|_{I\setminus\{h\}}).
\]
Strong induction on $\sum_i a_i+|I|$, starting from the empty index set,
proves \eqref{eq:good}.
\end{proof}

Take $I=\{1,\ldots,n\}$ and $a_i=q$. On the unit circle,
\[
\prod_{i\ne j}\left(1-\frac{z_i}{z_j}\right)^q=Y(t)^q.
\]
The torus integral of a finite Laurent polynomial is its constant term,
so
\begin{equation}\label{eq:Y-integer-moments}
\int_{\T^n}Y(t)^q\dd t=\frac{(nq)!}{(q!)^n}
\qquad(q\in\mathbb Z_{\ge0}).
\end{equation}

\subsection{A Beta Product and the Gauss Multiplication Formula}

For the moment assume $n\ge2$. Let $B_1,\ldots,B_{n-1}$ be mutually
independent random variables with
\[
B_r\sim\operatorname{Beta}\left(\frac rn,1-\frac rn\right)
\qquad(1\le r\le n-1),
\]
and set
\begin{equation}\label{eq:Z}
 Z=n^n\prod_{r=1}^{n-1}B_r.
\end{equation}

\begin{lemma}[Mellin transform of the Beta product]
\label{lem:beta-mellin}
Assume $n\ge2$. For every $s\ge0$,
\begin{equation}\label{eq:Z-mellin}
\E[Z^s]=\frac{\Gamma(1+ns)}{\Gamma(1+s)^n}.
\end{equation}
\end{lemma}

\begin{proof}
The Beta integral, independence, and
$B(x,y)=\Gamma(x)\Gamma(y)/\Gamma(x+y)$ give
\begin{align}
\E[Z^s]
&=n^{ns}\prod_{r=1}^{n-1}
\frac{B(s+r/n,1-r/n)}{B(r/n,1-r/n)}\notag\\
&=\frac{n^{ns}}{\Gamma(1+s)^{n-1}}
\prod_{r=1}^{n-1}
\frac{\Gamma(s+r/n)}{\Gamma(r/n)}.
\label{eq:beta-before-gauss}
\end{align}
For completeness, we do not invoke the Gauss formula as an unexplained
black box. For $x>0$, define
\[
G_n(x)=n^{x-1}
\frac{\prod_{r=0}^{n-1}\Gamma((x+r)/n)}
{\prod_{r=0}^{n-1}\Gamma((1+r)/n)}.
\]
Using $\Gamma(u+1)=u\Gamma(u)$ and cancelling consecutive factors gives
\[
 G_n(x+1)=xG_n(x),\qquad G_n(1)=1.
\]
Clearly $G_n(x)>0$, and
\[
\log G_n(x)
=(x-1)\log n+\sum_{r=0}^{n-1}\log\Gamma((x+r)/n)-\text{constant}
\]
is convex. The Bohr--Mollerup uniqueness theorem
\cite[pp.~34--36]{Andrews1999} therefore yields
$G_n(x)=\Gamma(x)$, or equivalently the normalized Gauss formula
\begin{equation}\label{eq:gauss-normalized}
\frac{\prod_{r=0}^{n-1}\Gamma((x+r)/n)}
{\prod_{r=0}^{n-1}\Gamma((1+r)/n)}
=n^{1-x}\Gamma(x).
\end{equation}
Set $x=1+ns$ in \eqref{eq:gauss-normalized} and reindex the numerator
and denominator products by $r=1,\ldots,n$. Isolating the factor
$\Gamma(1+s)/\Gamma(1)$ corresponding to $r=n$ then gives
\[
n^{ns}\prod_{r=1}^{n-1}
\frac{\Gamma(s+r/n)}{\Gamma(r/n)}
=\frac{\Gamma(1+ns)}{\Gamma(1+s)}.
\]
Substitution into \eqref{eq:beta-before-gauss} proves
\eqref{eq:Z-mellin}.
\end{proof}

\subsection{Pushforward Measures and Moment Determinacy on a Compact Interval}

\begin{lemma}[Moment determinacy on a compact interval]
\label{lem:moment-unique}
Let $\mu,\nu$ be probability measures on the same compact interval
$[0,C]$. If
\[
 \int_0^C x^q\dd\mu(x)=\int_0^C x^q\dd\nu(x)
 \qquad(q\in\mathbb Z_{\ge0}),
\]
then, for every $f\in C([0,C])$,
\[
 \int_0^C f\dd\mu=\int_0^C f\dd\nu.
\]
\end{lemma}

\begin{proof}
Equality of the moments implies by linearity that the integrals of all
polynomials agree. Given $f\in C([0,C])$ and $\varepsilon>0$, the
Weierstrass approximation theorem \cite{Rudin1976} provides a polynomial
$p$ such that $\|f-p\|_\infty<\varepsilon$. Since both probability
measures have total mass $1$,
\[
\left|\int(f-p)\dd\mu\right|\le\varepsilon,\qquad
\left|\int(f-p)\dd\nu\right|\le\varepsilon.
\]
The integrals of $p$ agree; letting $\varepsilon\downarrow0$ proves the
claim.
\end{proof}

\begin{proposition}[The circular Dyson integral at real parameters]
\label{prop:dyson-real}
Assume $n\ge2$. For $s\ge0$,
\begin{equation}\label{eq:dyson-real}
\int_{\T^n}Y(t)^s\dd t
=\frac{\Gamma(1+ns)}{\Gamma(1+s)^n}.
\end{equation}
\end{proposition}

\begin{proof}
The pairwise estimate $|z_i-z_j|\le2$ gives
\[
 0\le Y\le4^{\binom n2}.
\]
Since $0\le B_r\le1$, we also have $0\le Z\le n^n$. Put
\[
C=4^{\binom n2}+n^n
\,.
\]
Let $\dd t^{\otimes n}$ denote the normalized product Haar probability
measure on $\T^n$, and let $\mathbb P$ be the probability measure
governing the mutually independent variables $B_1,\ldots,B_{n-1}$.
Define on $[0,C]$ the pushforward probability measures
\[
\mu=Y_*(\dd t^{\otimes n}),\qquad
\nu=Z_*(\dd\mathbb P).
\]
For every integer $q\ge0$, \eqref{eq:Y-integer-moments} and
Lemma~\ref{lem:beta-mellin} give
\[
\int x^q\dd\mu
=\frac{(nq)!}{(q!)^n}
=\int x^q\dd\nu.
\]
For $s>0$, Lemma~\ref{lem:moment-unique} applies to the continuous
function $f(x)=x^s$; for $s=0$, take $f\equiv1$. Therefore
\[
\int_{\T^n}Y(t)^s\dd t=\E[Z^s].
\]
Equation \eqref{eq:Z-mellin} now proves \eqref{eq:dyson-real}.
\end{proof}

\subsection{The Dyson Base Case}

Set $s=K/2=k+1/2$. Propositions~\ref{prop:circularization} and
\ref{prop:dyson-real} give
\begin{equation}\label{eq:base-gamma}
L(0,0)=\frac{m!\pi^m}{n!}
\frac{\Gamma(1+nK/2)}{\Gamma(1+K/2)^n}.
\end{equation}
Since both $n$ and $K$ are odd,
\[
\Gamma\left(1+\frac r2\right)
=\frac{r!!\sqrt\pi}{2^{(r+1)/2}}
\qquad(r\ \text{odd}),
\]
and hence
\[
\frac{\Gamma(1+nK/2)}{\Gamma(1+K/2)^n}
=\frac{2^m}{\pi^m}\frac{(nK)!!}{(K!!)^{n}}.
\]
Moreover,
$n!=n!!(2m)!!=n!!\,2^m m!$, so
\begin{equation}\label{eq:base}
 L(0,0)=\frac1{n!!}\frac{(nK)!!}{(K!!)^{n}}.
\end{equation}
At $a=b=0$, the right-hand side of Theorem~\ref{thm:main} telescopes:
\[
\frac1{n!!}\prod_{i=0}^{n-1}
\frac{((i+1)K)!!}{(iK)!!K!!}
=\frac1{n!!}\frac{(nK)!!}{(K!!)^{n}}.
\]
Thus the base case has been established unconditionally.

\section{The Laurent--Euler Recurrence in the Endpoint Parameters}\label{sec:shift}

Set
\begin{equation}\label{eq:Er}
E_r(X)=\sum_{\substack{T\subseteq\{1,\ldots,n\}\\|T|=r}}
\prod_{j\in T}(-x_j)=(-1)^re_r(X),
\end{equation}
and define
\[
H_r(a,b)=\Lambda(\Phi_{a,b,k}E_r),\qquad 0\le r\le n.
\]
In particular, $H_0(a,b)=L(a,b)$.  To treat boundary indices
uniformly, for $J\subseteq\{1,\ldots,n\}$ and $s\in\mathbb Z$ define
\[
E_s^{(J)}=
\sum_{\substack{T\subseteq\{1,\ldots,n\}\setminus J\\|T|=s}}
\prod_{j\in T}(-x_j),
\]
with the convention that $E_s^{(J)}=0$ when $s<0$ or
$s>n-|J|$.  We abbreviate
$E_s^{(i)}=E_s^{(\{i\})}$ and
$E_s^{(i,j)}=E_s^{(\{i,j\})}$.  For a fixed $i$, put
\[
P_{r,i}=(1-x_i)E_r^{(i)},\qquad
D_i=x_i\frac{\partial}{\partial x_i}.
\]

\subsection{A Laurent-Polynomial Identity}

\begin{lemma}[Local shift identity]\label{lem:local-shift}
If $0\le r<n$, then
\begin{align}
2\sum_{i=1}^{n}D_i(P_{r,i}\Phi_{a,b,k})
={}&-(n-r)(2b+rK)\Phi_{a,b,k}E_r\notag\\
&+(r+1)\bigl(2a+2+(n-r-1)K\bigr)
\Phi_{a,b,k}E_{r+1}.
\label{eq:local-shift}
\end{align}
\end{lemma}

\begin{proof}
We first establish the intermediate identities used below.  Elementary
counting of finite subsets gives
\begin{equation}\label{eq:subset-count}
\sum_iE_r^{(i)}=(n-r)E_r,\qquad
\sum_i(-x_i)E_r^{(i)}=(r+1)E_{r+1}.
\end{equation}
For $i\ne j$, define
\[
Q_{r,ij}=(1-x_i-x_j)E_r^{(i,j)}
+x_ix_jE_{r-1}^{(i,j)}.
\]
Decomposing subsets according to whether they contain $i$ and $j$
immediately yields
\begin{equation}\label{eq:P-difference}
x_iP_{r,i}-x_jP_{r,j}=(x_i-x_j)Q_{r,ij}.
\end{equation}
We next count the multiplicity of each fixed monomial.  A given
$r$-element subset occurring in $E_r$ appears in $E_r^{(i,j)}$
$\binom{n-r}{2}$ times.  A given $(r+1)$-element subset occurring in
$E_{r+1}$ appears $(r+1)(n-r-1)$ times in
$-x_iE_r^{(i,j)}-x_jE_r^{(i,j)}$ and $\binom{r+1}{2}$ times in
$x_ix_jE_{r-1}^{(i,j)}$.  Therefore
\begin{align}
2\sum_{i<j}Q_{r,ij}
={}&(n-r-1)(n-r)E_r\notag\\
&+(r+1)\bigl(2(n-r-1)+r\bigr)E_{r+1}.
\label{eq:Q-count}
\end{align}

For integers $d,t,q\ge0$, set
\[
\Psi_{d,t,q}=\Delta^{q+1}\prod_i x_i^{-d}(1-x_i)^t.
\]
In the field of Laurent rational functions,
\[
\frac{D_i\Psi_{d,t,q}}{\Psi_{d,t,q}}
=-d-\frac{t x_i}{1-x_i}
+(q+1)\sum_{j\ne i}\frac{x_i}{x_i-x_j}.
\]
Since $D_iP_{r,i}=-x_iE_r^{(i)}$, the Leibniz rule and
\eqref{eq:subset-count} show that the endpoint contribution is
\[
-2d(n-r)\Psi E_r
+2(r+1)(t+1-d)\Psi E_{r+1}.
\]
Combining the Vandermonde contribution in unordered pairs and applying
\eqref{eq:P-difference}, we obtain
\[
\begin{aligned}
2(q+1)\Psi\sum_{i<j}
\frac{x_iP_{r,i}-x_jP_{r,j}}{x_i-x_j}
&=2(q+1)\Psi\sum_{i<j}Q_{r,ij}.
\end{aligned}
\]
Substitution of \eqref{eq:Q-count} now gives
\begin{align}
2\sum_iD_i(P_{r,i}\Psi_{d,t,q})
={}&\bigl[-2d(n-r)+(q+1)(n-r-1)(n-r)\bigr]\Psi E_r
\notag\\
&+\bigl[2(r+1)(t+1-d)\notag\\
&\qquad +(q+1)(r+1)(2(n-r-1)+r)\bigr]\Psi E_{r+1}.
\label{eq:general-euler}
\end{align}
The calculation was performed in the field of Laurent rational
functions.  Nevertheless, the left-hand side of
\eqref{eq:general-euler} is a sum of Euler derivatives of Laurent
polynomials, while $\Psi E_r$ and $\Psi E_{r+1}$ on the right-hand
side are Laurent polynomials.  Thus \eqref{eq:general-euler} is also
an identity in the Laurent-polynomial ring: all formal denominators
$1-x_i$ and $x_i-x_j$ have disappeared from the final expression.

Finally use the normalized formula \eqref{eq:normalized-kernel}, that
is, substitute
\[
d=b+mK,\qquad t=a+b,\qquad q+1=K,\qquad n=2m+1.
\]
The coefficient of $E_r$ simplifies to
$-(n-r)(2b+rK)$, and that of $E_{r+1}$ to
$(r+1)(2a+2+(n-r-1)K)$.  The constant $\varepsilon$ is unaffected by
Euler differentiation, proving \eqref{eq:local-shift}.
\end{proof}

\subsection{Finite Logarithmic Weight Shifts: All Exponent Branches}

Fix the $q$th logarithmic pair
\[
u=2q-1,\qquad v=2q,
\]
and let $w^{[q]}$ be the weight obtained by deleting this logarithmic
pair.  Thus it requires
\[
\alpha_u=\alpha_v=0,\qquad \alpha_n=0,
\]
each remaining pair must still satisfy \eqref{eq:weight}, and its
weight is the product of the remaining $m-1$ factors.  Denote the
corresponding finite functional by $\Lambda^{[q]}$.

\begin{lemma}[Weight-shift identity]\label{lem:weight-scale}
If $\alpha'=\alpha+\bm e_u-\bm e_v$, then, for every
$\alpha\in\mathbb Z^n$,
\begin{equation}\label{eq:weight-scale}
\alpha_u w(\alpha)-(\alpha_u+1)w(\alpha')
=w^{[q]}(\alpha).
\end{equation}
\end{lemma}

\begin{proof}
First suppose that all pairs other than $(u,v)$ and the unpaired
coordinate satisfy their respective conditions, and write $W$ for the
product of their weights.

If $\alpha_u=p>0$ and $\alpha_v=-p$, then
\[
w(\alpha)=-W/p,\qquad w(\alpha')=-W/(p+1),
\]
so the left-hand side is $-W+W=0$; the right-hand side also vanishes
because $(\alpha_u,\alpha_v)\ne(0,0)$.  If
$(\alpha_u,\alpha_v)=(0,0)$, then $w(\alpha)=0$ and
$w(\alpha')=-W$, so the left-hand side is $W$, exactly
$w^{[q]}(\alpha)$.

It remains to check every other case.  If $\alpha_u<0$, neither
$\alpha_u$ nor $\alpha_u+1$ is a positive left exponent; in the sole
exception that needs to be separated, $\alpha_u=-1$, the second
coefficient $\alpha_u+1$ itself is zero.  Thus both terms still vanish.
If $\alpha_u=0$ but $\alpha_v\ne0$, then $\alpha$ is inadmissible,
whereas admissibility of $\alpha'$ would force
$\alpha_v-1=-1$, hence $\alpha_v=0$, a contradiction.  If
$\alpha_u>0$ but $\alpha_v\ne-\alpha_u$, then $\alpha$ is
inadmissible; admissibility of $\alpha'$ would similarly force
$\alpha_v-1=-(\alpha_u+1)$, hence
$\alpha_v=-\alpha_u$, again a contradiction.  Finally, if any other
pair is mismatched or $\alpha_n\ne0$, all three quantities
$w(\alpha)$, $w(\alpha')$, and $w^{[q]}(\alpha)$ vanish.  This exhausts
all exponent branches.
\end{proof}

\begin{corollary}[Euler boundary formula]\label{cor:euler-boundary}
For arbitrary finite Laurent polynomials $G$ and $F$,
\begin{align}
\Lambda\!\left(D_u((x_v-x_u)G)\right)
&=\Lambda^{[q]}(x_vG),\label{eq:euler-boundary}\\
\Lambda(D_vF)&=-\Lambda(D_uF).
\label{eq:paired-euler}
\end{align}
\end{corollary}

\begin{proof}
Write $G=\sum_\beta c_\beta X^\beta$.  The contribution of
$X^\beta$ to the left-hand side of the first identity is
\[
c_\beta\bigl[
\beta_u w(\beta+\bm e_v)
-(\beta_u+1)w(\beta+\bm e_u)
\bigr].
\]
Taking $\alpha=\beta+\bm e_v$ in Lemma~\ref{lem:weight-scale}, this is
exactly $c_\beta w^{[q]}(\beta+\bm e_v)$, the corresponding
contribution on the right.  Finite summation proves
\eqref{eq:euler-boundary}.  The second identity follows monomial by
monomial: whenever $w(\alpha)\ne0$, one has
$\alpha_v=-\alpha_u$, while inadmissible monomials contribute zero to
both sides.
\end{proof}

\subsection{Invariance of the Omitted-Log Functional and Boundary Vanishing}

Let $w=n$ denote the unpaired index.  Because $\Lambda^{[q]}$ requires
the exponents of $u,v,w$ all to be zero, it is invariant under every
permutation of these three coordinates.  For $\tau\in S_n$, we use the
convention
\[
(F\circ\tau)(x_1,\ldots,x_n)
:=F(x_{\tau(1)},\ldots,x_{\tau(n)}).
\]

\begin{lemma}[Permutation invariance of the omitted-log functional]
\label{lem:omitted-invariance}
If $\tau$ permutes only $\{u,v,w\}$, then, for every finite Laurent
polynomial $F$,
\begin{equation}\label{eq:omitted-invariance}
\Lambda^{[q]}(F\circ\tau)=\Lambda^{[q]}(F).
\end{equation}
\end{lemma}

\begin{proof}
In $F=\sum_\alpha c_\alpha X^\alpha$, reindex the coefficients by
$\alpha\mapsto\alpha\circ\tau^{-1}$.  Whenever $w^{[q]}$ is nonzero,
the three permuted coordinates must all be zero, whereas the other
paired coordinates are unchanged.  Consequently
$w^{[q]}(\alpha\circ\tau^{-1})=w^{[q]}(\alpha)$, and the finite sum
\eqref{eq:Lambda} is invariant.
\end{proof}

Decomposition by subsets gives
\begin{equation}\label{eq:P-remainder}
P_{r,u}-P_{r,v}=(x_v-x_u)R_{r,uv},
\end{equation}
where
\[
R_{0,uv}=1,\qquad
R_{r,uv}=E_r^{(u,v)}-E_{r-1}^{(u,v)}\quad(r>0).
\]

\begin{lemma}[Three-variable boundary vanishing]\label{lem:boundary-zero}
\begin{equation}\label{eq:boundary-zero}
\Lambda^{[q]}(x_vR_{r,uv}\Phi_{a,b,k})=0.
\end{equation}
\end{lemma}

\begin{proof}
By \eqref{eq:normalized-kernel} and the oddness of $K$, interchanging
any two variables changes the sign of $\Phi_{a,b,k}$.

Consider first a subset $T$ in the expansion of $E_d^{(u,v)}$, and the
corresponding monomial
\[
M_T=x_v\prod_{j\in T}(-x_j)\Phi_{a,b,k},
\qquad T\cap\{u,v\}=\varnothing.
\]
If $w\notin T$, interchange $u$ and $w$.  Since neither $u$ nor $w$
lies in $T$, the factor $x_v\prod_{j\in T}(-x_j)$ is fixed;
$\Phi$ changes sign; and Lemma~\ref{lem:omitted-invariance} says that
the functional is unchanged.  Hence
\[
\Lambda^{[q]}(M_T)=-\Lambda^{[q]}(M_T),
\]
and this value is zero.

If $w\in T$, write
\[
M_T=x_v(-x_w)\prod_{j\in T\setminus\{w\}}(-x_j)\Phi.
\]
Interchange $v$ and $w$.  The first two factors change from
$x_v(-x_w)$ to $x_w(-x_v)$ and hence their product is fixed; the
remaining subset avoids $v,w$; $\Phi$ again changes sign; and the
functional is again invariant by Lemma~\ref{lem:omitted-invariance}.
This term is therefore also zero.  The argument covers every subset in
$E_d^{(u,v)}$, including the empty subset when $d=0$.  Since
$R_{r,uv}$ is a difference of two such polynomials,
\eqref{eq:boundary-zero} follows.
\end{proof}

\begin{proposition}[Endpoint-shift recurrence]\label{prop:shift-recurrence}
For $0\le r<n$,
\begin{equation}\label{eq:shift-recurrence}
(n-r)(2b+rK)H_r(a,b)
=(r+1)\bigl(2a+2+(n-r-1)K\bigr)H_{r+1}(a,b).
\end{equation}
\end{proposition}

\begin{proof}
Pair $D_u(P_{r,u}\Phi)$ with $D_v(P_{r,v}\Phi)$.  By
\eqref{eq:paired-euler}, \eqref{eq:P-remainder}, and
\eqref{eq:euler-boundary},
\[
\begin{aligned}
&\Lambda\bigl(D_u(P_{r,u}\Phi)+D_v(P_{r,v}\Phi)\bigr)\\
&\quad=\Lambda\bigl(D_u((P_{r,u}-P_{r,v})\Phi)\bigr)
=\Lambda^{[q]}(x_vR_{r,uv}\Phi)=0,
\end{aligned}
\]
where the last equality is Lemma~\ref{lem:boundary-zero}.  Sum over all
$m$ pairs.  For the remaining unpaired coordinate $w=n$, every
admissible exponent is zero, so $\Lambda(D_wF)=0$.  Therefore
\begin{equation}\label{eq:total-euler-zero}
\Lambda\left(\sum_{i=1}^{n}D_i(P_{r,i}\Phi)\right)=0.
\end{equation}
Applying $\Lambda$ to Lemma~\ref{lem:local-shift}, using
\eqref{eq:total-euler-zero}, and rearranging proves
\eqref{eq:shift-recurrence}.
\end{proof}

\section{Chu--Vandermonde Summation and the Main Theorem}

The coefficient
\[
(r+1)\bigl(2a+2+(n-r-1)K\bigr)
\]
on the right-hand side of \eqref{eq:shift-recurrence} is strictly
positive.  Thus, starting from $H_0(a,b)=L(a,b)$, we solve the
recurrence by dividing only by this coefficient and never by
$2b+rK$.  Induction gives
\begin{equation}\label{eq:Hr}
H_r(a,b)=\binom nr
\frac{\prod_{j=0}^{r-1}(2b+jK)}
{\prod_{j=0}^{r-1}(2a+2+(n-1-j)K)}
L(a,b).
\end{equation}
Empty products are understood to be $1$.  In particular, when $b=0$
and $r=0$, the recurrence first gives $H_1(a,0)=0$; induction then
gives $H_r(a,0)=0$ for every $r>0$.  This agrees exactly with the zero
factor at $j=0$ in the numerator of \eqref{eq:Hr}, so the boundary case
involves no division by zero.

\begin{lemma}[Finite-step Chu--Vandermonde identity]\label{lem:vandermonde}
If all denominators below are nonzero, then
\begin{equation}\label{eq:vandermonde-ratio}
\sum_{r=0}^{n}\binom nr
\frac{\prod_{j=0}^{r-1}(B+jK)}
{\prod_{j=0}^{r-1}(A+(n-1-j)K)}
=\prod_{i=0}^{n-1}\frac{A+B+iK}{A+iK}.
\end{equation}
\end{lemma}

\begin{proof}
We first prove the denominator-free polynomial identity
\begin{equation}\label{eq:vandermonde-poly}
\sum_{r=0}^{n}\binom nr
\prod_{j=0}^{r-1}(B+jK)
\prod_{j=0}^{n-r-1}(A+jK)
=\prod_{i=0}^{n-1}(A+B+iK).
\end{equation}
For $n=0$, both sides are $1$.  In the induction step from $n$ to
$n+1$, apply Pascal's identity and shift the index in the second
part from $r$ to $r-1$.  For each fixed $0\le r\le n$, the two
coefficients combine as
\[
\bigl(A+(n-r)K\bigr)+\bigl(B+rK\bigr)=A+B+nK.
\]
The left-hand side is therefore $A+B+nK$ times its counterpart of
order $n$, and the induction hypothesis proves
\eqref{eq:vandermonde-poly}.  Finally multiply the $r$th summand of
\eqref{eq:vandermonde-ratio} by
$\prod_{i=0}^{n-1}(A+iK)$.  Cancelling the last $r$ factors produces
precisely the $r$th summand of \eqref{eq:vandermonde-poly}.  Division
by the nonzero product completes the proof.
\end{proof}

Since
\[
\sum_{r=0}^{n}E_r(X)=\prod_{i=1}^{n}(1-x_i),
\]
finite linearity gives
\begin{equation}\label{eq:H-sum}
\sum_{r=0}^{n}H_r(a,b)=L(a+1,b).
\end{equation}
Take
\[
A=2a+2,\qquad B=2b
\]
in Lemma~\ref{lem:vandermonde}; all denominators are positive.
Substitution of \eqref{eq:Hr} into \eqref{eq:H-sum} yields
\begin{equation}\label{eq:L-shift}
L(a+1,b)=L(a,b)
\prod_{i=0}^{n-1}
\frac{2a+2b+2+iK}{2a+2+iK}.
\end{equation}

Let $R(a,b)$ denote the explicit right-hand side of
Theorem~\ref{thm:main}.  Applying $(N+2)!!=(N+2)N!!$ factor by factor
gives
\begin{equation}\label{eq:R-shift}
R(a+1,b)=R(a,b)
\prod_{i=0}^{n-1}
\frac{2a+2b+2+iK}{2a+2+iK}.
\end{equation}
Clearly $R(a,b)=R(b,a)$, while Corollary~\ref{cor:swap} gives
$L(a,b)=L(b,a)$.

\begin{proof}[Proof of Theorem~\ref{thm:main}]
Section~\ref{sec:dyson} established $L(0,0)=R(0,0)$.  First fix
$b=0$.  Induction on $a$ using \eqref{eq:L-shift} and
\eqref{eq:R-shift} gives $L(a,0)=R(a,0)$.  Symmetry in $a,b$ then gives
$L(0,b)=R(0,b)$.  Finally fix an arbitrary $b$ and induct once more on
$a$.  This proves $L(a,b)=R(a,b)$ for all nonnegative integers $a,b$,
which is precisely \eqref{eq:main}.
\end{proof}

\section{Proof Dependencies and Lean Semantic Audit}
\label{sec:lean-audit}

\subsection{Logical Closure}

The logical dependencies of the preceding proof are
\begin{align*}
\text{Good--Dyson}
&\Longrightarrow \{\mathbb E(Y^r):r\in\mathbb Z_{\ge0}\},\\
\text{beta product and compact moment determinacy}
&\Longrightarrow \mathbb E(Y^{K/2}),\\
\text{Fourier bridge and augmented Pfaffian}
&\Longrightarrow L(0,0)=R(0,0),\\
\text{weight shift and three-variable boundary vanishing}
&\Longrightarrow \text{endpoint-shift recurrence},\\
\text{Chu--Vandermonde and endpoint symmetry}
&\Longrightarrow L(a,b)=R(a,b).
\end{align*}
The first implication uses only finite Laurent identities with
nonnegative integer exponents.  The unique extension from integer to
half-integer moments takes place between probability measures supported
on the same compact interval; it does not use the complex Morris
conjecture.

\subsection{Correspondence Between the Mathematical Objects and the Lean Interfaces}

The final theorem in the formal development is
\[
\texttt{LogarithmicMorrisFull.logarithmicMorris\_full}.
\]
Its type proves \texttt{LogarithmicMorrisStatement S} for an arbitrary
\texttt{Setup}. The following table records correspondences at the level
of mathematical objects and formal interfaces; it does not assert that
every natural-language argument below has been formalized line by line.
The explicit definitions in Sections~1 and~2 provide the human-checkable
identification of the formal statement with the constant-term identity
of \cite[Theorem~1.3]{LogMorris}.

\begin{longtable}{>{\raggedright\arraybackslash}p{0.31\textwidth}
                  >{\raggedright\arraybackslash}p{0.61\textwidth}}
\toprule
Mathematical object & Concrete Lean interface or theorem\\
\midrule
\endfirsthead
\toprule
Mathematical object & Concrete Lean interface or theorem\\
\midrule
\endhead
$n=2m+1,\ K=2k+1$ &
\texttt{Setup.odd\_rank} and \texttt{Setup.K}\\
The finite Laurent kernel in \eqref{eq:kernel} &
\texttt{morrisKernel}; its coefficient field is $\mathbb Q$, and its
exponent type is \texttt{Fin n -> Z}\\
The weight \eqref{eq:weight} and functional \eqref{eq:Lambda} &
\texttt{StandardLogAdmissible}, \texttt{standardLogWeight}, and
\texttt{standardLogCT}\\
The two sides of the theorem &
\texttt{logarithmicMorrisLHS} and \texttt{logarithmicMorrisRHS}\\
The Fourier bridge and $L^1$ control &
\texttt{ScratchCircleLogFourier},
\texttt{ScratchPairedLogIntegrable}, and
\texttt{ScratchLogCTBridge}\\
The alternating sum of augmented sawtooth values &
\texttt{LogarithmicMorrisPfaffian}; Lean proves the same finite
alternating-sum identity by a wholly finite permutation-sum argument,
rather than formalizing line by line the augmented-matrix congruence
argument in this section\\
Ordered chambers, arbitrary permutations, and collisions &
\texttt{ScratchLogPfaff}, \texttt{ScratchPointwise}, and
\texttt{ScratchPermutationIntegral}\\
Good--Dyson integer moments &
\texttt{ScratchGoodBoundary} and
\texttt{LogarithmicMorrisEvenMoments}\\
The beta product and Gauss's formula &
\texttt{LogarithmicMorrisDysonBetaProduct} and
\texttt{AristotleGammaMultiplicationFresh}\\
Pushforward measures and compactly supported moment determinacy &
\texttt{LogarithmicMorrisDysonMomentComparison} and
\texttt{AristotleMomentUniquenessFresh}\\
The local Euler identity &
\texttt{ScratchShiftLocal}\\
All branches of the weight-shift identity and three-variable boundary
vanishing &
\texttt{ScratchShiftLogZero}\\
The parameter recurrence and final induction &
\texttt{LogarithmicMorrisParameterReduction} and
\texttt{LogarithmicMorrisFull}\\
\bottomrule
\end{longtable}

\begin{remark}[The circular logarithm at collision points]
Lean's \texttt{Complex.log} is a total function, with value $0$ at
$0$.  The formal files may therefore state identities that hold at
collision points as well: the Vandermonde factor makes the Laurent
kernel vanish there.  In ordinary mathematical usage, the principal
logarithm is generally regarded as defined only away from zero.  We
accordingly used the extension $\Logc(0)=0$ and performed the analytic
calculation in Proposition~\ref{prop:circularization} only on the
full-measure set of pairwise distinct coordinates.  The two conventions
give the same integrals and Fourier coefficients.
\end{remark}

\begin{remark}[The rank-one case]
When $m=0$, the logarithmic product over \texttt{Fin m} is empty, and
\texttt{StandardLogAdmissible} requires only that the exponent of the
unique coordinate be zero.  Thus \texttt{standardLogCT} genuinely
reduces to the ordinary constant term.  The file proving the circular
Dyson evaluation also treats $n=1$ as a separate branch, in agreement
with the direct calculation in Section~1.
\end{remark}

\subsection{Reproducible Build and Axiom Report}

The complete release repository is publicly available at
\begin{center}
\url{https://github.com/halfblood12312300-cell/logarithmic-morris-lean}.
\end{center}
The Lean proof snapshot audited in this paper is the immutable commit
\begin{center}
\texttt{321df3a4af7231df43bd5b976ac366e44d0a1ce6}.
\end{center}
At this audited snapshot, the repository root contains \texttt{lakefile.toml},
\texttt{lake-manifest.json}, and \texttt{lean-toolchain}.  The last of
these pins
\[
\texttt{leanprover/lean4:v4.28.0},
\]
and that Lake configuration pins Mathlib at \texttt{v4.28.0}.  On
22 July 2026, the following commands were run from a fresh checkout:
\begin{verbatim}
git clone https://github.com/halfblood12312300-cell/logarithmic-morris-lean.git
cd logarithmic-morris-lean
git checkout 321df3a4af7231df43bd5b976ac366e44d0a1ce6
lake exe cache get
lake build logarithmicmorris.LogarithmicMorrisFull
lake env lean --trust=0 logarithmicmorris/LogarithmicMorrisAudit.lean
\end{verbatim}
The final-target build succeeded (8082 tasks).  The final audit command
produced
the following output:
\begin{verbatim}
LogarithmicMorrisFull.logarithmicMorris_full
  (S : LogarithmicMorrisFull.Setup) :
  LogarithmicMorrisFull.LogarithmicMorrisStatement S
'LogarithmicMorrisFull.logarithmicMorris_full' depends on axioms:
  [propext, Classical.choice, Quot.sound]
\end{verbatim}
Thus the axiom closure contains no \texttt{sorryAx} and no external
axiom introduced for this problem.  The three reported axioms are
Lean/Mathlib's standard logical dependencies for quotient types,
propositional extensionality, and classical choice; none encodes a
Morris, Dyson, Selberg, or integral-evaluation assumption.  The project
release contains exactly the local transitive import closure required by
the final theorem, together with its public entry point and audit file.
No Lean source file in this release contains a \texttt{sorry}
declaration.  The absence of an unreported logical assumption is not
inferred from file names or a textual search: it is certified by the
kernel-level \texttt{\#print axioms} report reproduced above.  The
repository also includes a GitHub Actions workflow that rebuilds the
complete release library, rejects \texttt{sorry}, \texttt{sorryAx}, and
local \texttt{axiom} declarations, runs the Lean~4.28
\texttt{leanchecker} environment verifier, and prints the same axiom
closure.

\section*{Relation to the Original Conditional Proof}

The unconditional route in this paper differs from the derivation in
the original article, which was conditional on the complex Morris
identity.  The conjectural input is replaced precisely in
Section~\ref{sec:dyson}: Good's integer Dyson identity determines all
moments of a compactly supported random variable, and a beta product
then evaluates its $K/2$ moment.  The formal verification ensures that
all finite algebraic branches, measure pushforwards, integrability
arguments, and parameter recurrences are closed.  The present paper
expands those machine-checked interfaces into a self-contained
mathematical proof.

\begingroup
\small
\begin{thebibliography}{9}
\bibitem{Dyson1962}
F. J. Dyson,
\newblock Statistical theory of the energy levels of complex systems I,
\newblock \emph{Journal of Mathematical Physics} 3 (1962), 140--156.

\bibitem{Good1970}
I. J. Good,
\newblock Short proof of a conjecture by Dyson,
\newblock \emph{Journal of Mathematical Physics} 11 (1970), 1884.

\bibitem{LogMorris}
T. Chappell, A. Lascoux, S. O. Warnaar, and W. Zudilin,
\newblock Logarithmic and complex constant term identities,
\newblock in \emph{Computational and Analytical Mathematics},
Springer Proceedings in Mathematics and Statistics 50
(2013), 219--250; arXiv:1112.3130.

\bibitem{Andrews1999}
G. E. Andrews, R. Askey, and R. Roy,
\newblock \emph{Special Functions},
\newblock Encyclopedia of Mathematics and its Applications 71,
Cambridge University Press, Cambridge, 1999.

\bibitem{Rudin1976}
W. Rudin,
\newblock \emph{Principles of Mathematical Analysis}, 3rd ed.,
\newblock McGraw--Hill, New York, 1976.

\bibitem{Knuth1996}
D. E. Knuth,
\newblock Overlapping Pfaffians,
\newblock \emph{Electronic Journal of Combinatorics} 3 (1996),
no.~2, Research Paper 5.
\end{thebibliography}
\endgroup

\end{document}
